\chapter{Sistema de Estimación Social (SSP)}
\label{anexo:sistema-estimacion-social}

\section{Sistema de estimación de ítems sociales (\glspl{social-story-points})}

La Sección~\ref{sec:sistema-estimacion-ssp} del Capítulo~\ref{chap:capitulo4} presenta la escala SSP, la fórmula de cálculo y la tabla de referencia rápida. Este anexo complementa esa presentación con la metodología detallada de estimación en 3~pasos, la matriz completa de severidad y las mejores prácticas de calibración.

\subsection{Escala de estimación SSP}

La escala SSP utiliza los mismos valores Fibonacci: 1, 2, 3, 5, 8 y 13. Valores superiores a 13 indican que el ítem debe descomponerse.

\begingroup
\scriptsize
\renewcommand{\arraystretch}{1.4}
\setlength{\tabcolsep}{4pt}
\setlength{\LTleft}{0pt}
\setlength{\LTright}{0pt}

\begin{longtable}{|
    >{\Centering\hspace{0pt}}p{1.5cm}|     % Col 1: SSP
    >{\Centering\hspace{0pt}}p{2.5cm}|     % Col 2: Duración
    >{\RaggedRight\hspace{0pt}}p{6.4cm}|   % Col 3: Tipo de intervención
    >{\RaggedRight\hspace{0pt}}p{4.5cm}|   % Col 4: Ejemplo
    }

    % --- 1. ENCABEZADO PARA LA PRIMERA PÁGINA ---
    \caption{Escala de Social Story Points (SSP)} \label{tab:escala-SSP}                                                                \\
    \hline
    \textbf{SSP} & \textbf{Duración} & \textbf{Tipo de intervención}                        & \textbf{Ejemplo}                          \\
    \hline
    \endfirsthead

    % --- 2. ENCABEZADO PARA PÁGINAS SIGUIENTES ---
    \hline
    \multicolumn{4}{|c|}{{\bfseries \tablename\ \thetable{} -- Continuación}}                                                             \\
    \hline
    \textbf{SSP} & \textbf{Duración} & \textbf{Tipo de intervención}                        & \textbf{Ejemplo}                          \\
    \hline
    \endhead

    % --- 3. PIE PARA PÁGINAS INTERMEDIAS ---
    \hline
    \multicolumn{4}{|r|}{\textit{Continúa en la siguiente página...}}                                                                     \\
    \hline
    \endfoot

    % --- 4. PIE FINAL ---
    \hline
    \endlastfoot

    1--2         & 1--4 horas        & Conversación facilitada, ajuste menor                & \textit{Feedback} rápido, norma simple    \\
    \hline
    3--5         & 4 h--2 días       & \textit{Workshop}, mediación, retrospectiva profunda & Código de conducta, mediación bilateral   \\
    \hline
    8--13        & 2 días--1 semana  & Intervención organizacional, crisis                  & Cambio de liderazgo, \gls{burnout} severo \\
\end{longtable}
\endgroup

\subsection{Metodología de estimación en 3 pasos}

La estimación SSP sigue un proceso de tres pasos: identificar el \gls{community-smells}, evaluar severidad, y evaluar complejidad de intervención.

\begingroup
\scriptsize
\renewcommand{\arraystretch}{1.4}
\setlength{\tabcolsep}{4pt}
\setlength{\LTleft}{0pt}
\setlength{\LTright}{0pt}

\begin{longtable}{|
    >{\Centering\hspace{0pt}}p{1.0cm}|     % Col 1: Paso
    >{\RaggedRight\hspace{0pt}}p{3.5cm}|   % Col 2: Criterio
    >{\RaggedRight\hspace{0pt}}p{5.9cm}|   % Col 3: Entrada
    >{\RaggedRight\hspace{0pt}}p{4.5cm}|   % Col 4: Salida
    }

    % --- 1. ENCABEZADO PARA LA PRIMERA PÁGINA ---
    \caption{Proceso de estimación de SSP} \label{tab:proceso-SSP}                                                          \\
    \hline
    \textbf{Paso} & \textbf{Criterio}          & \textbf{Entrada}                      & \textbf{Salida}                    \\
    \hline
    \endfirsthead

    % --- 2. ENCABEZADO PARA PÁGINAS SIGUIENTES ---
    \hline
    \multicolumn{4}{|c|}{{\bfseries \tablename\ \thetable{} -- Continuación}}                                                 \\
    \hline
    \textbf{Paso} & \textbf{Criterio}          & \textbf{Entrada}                      & \textbf{Salida}                    \\
    \hline
    \endhead

    % --- 3. PIE PARA PÁGINAS INTERMEDIAS ---
    \hline
    \multicolumn{4}{|r|}{\textit{Continúa en la siguiente página...}}                                                         \\
    \hline
    \endfoot

    % --- 4. PIE FINAL ---
    \hline
    \endlastfoot

    1             & Identificar \textit{smell} & Registro de \textit{Community Smells} & Tipo específico confirmado         \\
    \hline
    2             & Evaluar severidad          & \gls{health-score} y observación      & Leve / Moderada / Severa / Crítica \\
    \hline
    3             & Evaluar complejidad        & Alcance, facilitación y sesiones      & Baja (+0) / Media (+1) / Alta (+2) \\
\end{longtable}
\endgroup

\subsubsection{Matriz de severidad y SSP base}

\begingroup
\scriptsize
\renewcommand{\arraystretch}{1.4}
\setlength{\tabcolsep}{4pt}
\setlength{\LTleft}{0pt}
\setlength{\LTright}{0pt}

\begin{longtable}{|
    >{\RaggedRight\hspace{0pt}}p{2.5cm}|   % Col 1: Severidad
    >{\RaggedRight\hspace{0pt}}p{10.2cm}|   % Col 2: Indicadores
    >{\Centering\hspace{0pt}}p{2.5cm}|     % Col 3: SSP base
    }

    % --- 1. ENCABEZADO PARA LA PRIMERA PÁGINA ---
    \caption{Niveles de severidad y SSP base} \label{tab:severidad-SSP}                                     \\
    \hline
    \textbf{Severidad} & \textbf{Indicadores}                                           & \textbf{SSP base} \\
    \hline
    \endfirsthead

    % --- 2. ENCABEZADO PARA PÁGINAS SIGUIENTES ---
    \hline
    \multicolumn{3}{|c|}{{\bfseries \tablename\ \thetable{} -- Continuación}}                                 \\
    \hline
    \textbf{Severidad} & \textbf{Indicadores}                                           & \textbf{SSP base} \\
    \hline
    \endhead

    % --- 3. PIE PARA PÁGINAS INTERMEDIAS ---
    \hline
    \multicolumn{3}{|r|}{\textit{Continúa en la siguiente página...}}                                         \\
    \hline
    \endfoot

    % --- 4. PIE FINAL ---
    \hline
    \endlastfoot

    Leve               & Afecta 1--2 personas ocasionalmente; no impacta la entrega     & 1--2              \\
    \hline
    Moderada           & Afecta 3--5 personas; retrasos menores; tensión visible        & 3--5              \\
    \hline
    Severa             & Afecta al equipo entero; riesgo de rotación; calidad degradada & 5--8              \\
    \hline
    Crítica            & Equipo disfuncional; rotación inminente; proyecto en riesgo    & 8--13             \\
\end{longtable}
\endgroup

\subsubsection{Cálculo del SSP final}

\begin{equation}
    \text{SSP Final} = \text{SSP Base (severidad)} + \text{Ajuste (complejidad)}
\end{equation}

\textbf{Ejemplo}: Conflicto Dev A-B (Severidad Moderada = 3 SSP base) + Mediación especializada requerida (Complejidad Alta = +2) $\rightarrow$ SSP Final = 5.

\subsection{Tabla de referencia rápida}

\begingroup
\scriptsize
\renewcommand{\arraystretch}{1.4}
\setlength{\tabcolsep}{4pt}
\setlength{\LTleft}{0pt}
\setlength{\LTright}{0pt}

\begin{longtable}{|
    >{\RaggedRight\hspace{0pt}}p{4.0cm}|   % Col 1: Community smell
    >{\Centering\hspace{0pt}}p{2.5cm}|     % Col 2: Severidad típica
    >{\RaggedRight\hspace{0pt}}p{6.9cm}|   % Col 3: Intervención típica
    >{\Centering\hspace{0pt}}p{1.5cm}|     % Col 4: SSP
    }

    % --- 1. ENCABEZADO PARA LA PRIMERA PÁGINA ---
    \caption{Referencia rápida de estimación SSP por olor comunitario} \label{tab:referencia-SSP}                                \\
    \hline
    \textbf{\textit{Community smell}} & \textbf{Severidad típica} & \textbf{Intervención típica}                  & \textbf{SSP} \\
    \hline
    \endfirsthead

    % --- 2. ENCABEZADO PARA PÁGINAS SIGUIENTES ---
    \hline
    \multicolumn{4}{|c|}{{\bfseries \tablename\ \thetable{} -- Continuación}}                                                      \\
    \hline
    \textbf{\textit{Community smell}} & \textbf{Severidad típica} & \textbf{Intervención típica}                  & \textbf{SSP} \\
    \hline
    \endhead

    % --- 3. PIE PARA PÁGINAS INTERMEDIAS ---
    \hline
    \multicolumn{4}{|r|}{\textit{Continúa en la siguiente página...}}                                                              \\
    \hline
    \endfoot

    % --- 4. PIE FINAL ---
    \hline
    \endlastfoot

    \textit{Radio Silence}            & Moderada                  & \textit{Workshop} y ajuste de proceso         & 3--5         \\
    \hline
    \textit{Lone Wolf}                & Leve--Moderada            & Mediación y reintegración                     & 3--5         \\
    \hline
    Conflicto interpersonal leve      & Leve                      & Mediación bilateral simple                    & 3            \\
    \hline
    Conflicto interpersonal severo    & Severa                    & Mediación intensiva y seguimiento             & 8            \\
    \hline
    \textit{Burnout} leve / moderado  & Leve--Moderada            & Conversación y ajuste de carga                & 2--5         \\
    \hline
    \textit{Burnout} severo           & Severa--Crítica           & Intervención de crisis                        & 13           \\
    \hline
    Inseguridad psicológica           & Moderada--Severa          & \textit{Workshop} y retrospectiva facilitada  & 5--8         \\
    \hline
    \textit{Knowledge hoarding}       & Moderada                  & Documentación y transferencia de conocimiento & 3--5         \\
\end{longtable}
\endgroup

\subsection{Conversión SSP \texorpdfstring{$\leftrightarrow$}{<->} SP y gestión de capacidad}

El principio fundamental es: \textbf{1 SSP $\approx$ 1 SP}. Ambos tipos de ítems pueden sumarse para calcular la capacidad total del \textit{Sprint}.

\begingroup
\scriptsize
\renewcommand{\arraystretch}{1.4}
\setlength{\tabcolsep}{4pt}
\setlength{\LTleft}{0pt}
\setlength{\LTright}{0pt}

\begin{longtable}{|
    >{\RaggedRight\hspace{0pt}}p{3.0cm}|   % Col 1: Tipo
    >{\RaggedRight\hspace{0pt}}p{7.9cm}|   % Col 2: Ítems ejemplo
    >{\Centering\hspace{0pt}}p{2.0cm}|     % Col 3: Puntos
    >{\Centering\hspace{0pt}}p{2.0cm}|     % Col 4: Total
    }

    % --- 1. ENCABEZADO PARA LA PRIMERA PÁGINA ---
    \caption{Ejemplo de cálculo de capacidad combinada} \label{tab:capacidad-combinada}                                                                              \\
    \hline
    \textbf{Tipo}                          & \textbf{Ítems ejemplo}                                                           & \textbf{Puntos} & \textbf{Total}     \\
    \hline
    \endfirsthead

    % --- 2. ENCABEZADO PARA PÁGINAS SIGUIENTES ---
    \hline
    \multicolumn{4}{|c|}{{\bfseries \tablename\ \thetable{} -- Continuación}}                                                                                          \\
    \hline
    \textbf{Tipo}                          & \textbf{Ítems ejemplo}                                                           & \textbf{Puntos} & \textbf{Total}     \\
    \hline
    \endhead

    % --- 3. PIE PARA PÁGINAS INTERMEDIAS ---
    \hline
    \multicolumn{4}{|r|}{\textit{Continúa en la siguiente página...}}                                                                                                  \\
    \hline
    \endfoot

    % --- 4. PIE FINAL ---
    \hline
    \endlastfoot

    Técnico                                & \textit{Feature} A (13 SP) + \textit{Feature} B (8 SP) + \textit{Bug fix} (5 SP) & 26 SP           &                    \\
    \hline
    Social                                 & Mediación de conflicto (5 SSP) + \textit{daily} asíncrona (3 SSP)                & 8 SSP           &                    \\
    \hline
    \textbf{Capacidad del \textit{Sprint}} &                                                                                  &                 & \textbf{34 puntos} \\
\end{longtable}
\endgroup

\subsubsection{Mejores prácticas de estimación}

\begin{itemize}
    \item \textbf{Estimar de manera conservadora}: Añadir 20~\% de \textit{buffer} en situaciones inciertas.
    \item \textbf{Descomponer ítems grandes}: Si un ítem supera 8 SSP, buscar cómo dividirlo.
    \item \textbf{Registrar real vs. estimado}: Mantener historial para calibración.
    \item \textbf{Un ítem = un \textit{community smell} específico}: Evitar ítems vagos como ``mejorar comunicación''.
\end{itemize}